\label{int:chap:producto-nativo}
\index{producto nativo}

La extensión euclídea residuo--cociente de APP contiene dos aplicaciones
locales: la aplicación aditiva normaliza colisiones y la multiplicativa
produce productos parciales. En esta sección se construye la multiplicación
global de palabras usando exclusivamente ambas aplicaciones. La correspondencia con la multiplicación entera será una
consecuencia del entrelazador, no la definición de la operación.

\section{Coordenadas euclídeas de la célula local}

\begin{definition}[Aplicaciones aditiva y multiplicativa sobre un mismo soporte]
Sea $\Dnine=\{1,\ldots,9\}$. Para $a,b\in\Dnine$ se definen por división
euclídea con resto en $\Dnine$:
\begin{align}
 a+b&=9q^+(a,b)+R^+(a,b),\label{int:eq:local-plus}\\
 ab&=9q^\times(a,b)+R^\times(a,b).\label{int:eq:local-times}
\end{align}
En particular, el resto nulo se representa por $9$ y el cociente se reduce en
una unidad. La cuádrupla
\[
 \CellRQ(a,b)=
 (R^+(a,b),q^+(a,b);R^\times(a,b),q^\times(a,b))
\]
es la extensión aritmética local por residuo y cociente.
\end{definition}

La agregación por las tres clases módulo tres de la tabla visible
\(R^\times\) define una matriz distinta de las aplicaciones locales:
\[
(M_\Pi)_{rs}
=\frac19
\sum_{\substack{a\equiv r\;(3)\\b\equiv s\;(3)}}R^\times(a,b),
\qquad r,s\in\{1,2,0\},
\]
y, en ese orden de clases,
\[
M_\Pi=
\begin{pmatrix}4&5&6\\5&4&6\\6&6&9\end{pmatrix}.
\]
Así, \(q^+\) y \(q^\times\) son coordenadas de cociente de una celda,
mientras que \(M_\Pi\) es un operador agregado sobre clases residuales.

Las $81$ celdas satisfacen exactamente
\[
\begin{array}{c|rrrr}
 &\sum R&\sum q&\sum(9q+R)\\\hline
 +&405&45&810\\
 \times&459&174&2025
\end{array}
\]
y por tanto
\[
 2025-810=(459-405)+9(174-45)=54+9\cdot129=1215.
\]
El cociente $q^\times$ es necesario para la reconstrucción. Por ejemplo,
\[
 2\cdot5=9\cdot1+1.
\]
Si se conserva sólo $R^\times$, el producto $10$ se confunde con el valor
visible $1$.

\subsection{Cociclos de asociatividad y distributividad}

La coherencia del transductor puede decidirse ya en la célula. Para todo
$a,b,c\in\Dnine$ se verifican las tres identidades
\begin{align}
q^+(a,b)+q^+(R^+(a,b),c)
&=q^+(b,c)+q^+(a,R^+(b,c)),\label{int:eq:cocycle-plus}\\
q^\times(R^\times(a,b),c)+cq^\times(a,b)
&=q^\times(a,R^\times(b,c))+aq^\times(b,c),
\label{int:eq:cocycle-times}\\
q^\times(a,R^+(b,c))+aq^+(b,c)
&=q^+(R^\times(a,b),R^\times(a,c))\notag\\
&\quad+q^\times(a,b)+q^\times(a,c).
\label{int:eq:cocycle-distributive}
\end{align}

\begin{theorem}[Coherencia local de la extensión euclídea]
Las identidades \eqref{int:eq:cocycle-plus}--\eqref{int:eq:cocycle-distributive},
junto con las igualdades correspondientes de residuos, expresan
respectivamente asociatividad aditiva, asociatividad multiplicativa y
distributividad en la extensión residuo--cociente. Las tres familias tienen
cero discrepancias en los $9^3=729$ triples.
\end{theorem}

\begin{proof}
Al desarrollar $(a+b)+c$ mediante \eqref{int:eq:local-plus}, el coeficiente de
nueve es el lado izquierdo de \eqref{int:eq:cocycle-plus}; al desarrollar
$a+(b+c)$ se obtiene el derecho. La unicidad de la división con resto en
$\Dnine$ iguala tanto residuos como cocientes. El mismo argumento aplicado a
$(ab)c=a(bc)$ da \eqref{int:eq:cocycle-times}. Finalmente, las dos expansiones de
$a(b+c)=ab+ac$ producen \eqref{int:eq:cocycle-distributive}. La enumeración
directa de los \(729\) casos de cada identidad confirma el cálculo
algebraico.
\end{proof}

\section{Biyectividad de la numeración}

El símbolo $9$ no se identifica con el cero. En $\CellRQ$, $9$ representa un
cruce completo de la frontera nonádica. La sustitución por los dígitos
$0,\ldots,8$ eliminaría el cociente euclídeo que interviene en el
entrelazamiento.

\begin{definition}[Palabras canónicas]
Sea
\[
 \Wnine^+=\bigsqcup_{\ell\geq1}\Dnine^\ell,
 \qquad \Wnine=\{\varnothing\}\sqcup\Wnine^+.
\]
Las palabras se escriben desde el dígito menos significativo. Definimos
\[
 \operatorname{ev}_9(u_0,\ldots,u_{\ell-1})
 =\sum_{j=0}^{\ell-1}u_j9^j,
 \qquad \operatorname{ev}_9(\varnothing)=0.
\]
\end{definition}

\begin{theorem}[Unicidad de la forma normal]
La evaluación $\operatorname{ev}_9$ es una biyección de $\Wnine^+$ sobre
$\N_{\geq1}$, y de $\Wnine$ sobre $\N_0$.
\end{theorem}

\begin{proof}
Las palabras de longitud $\ell$ cubren el intervalo
\[
 \frac{9^\ell-1}{8}
 \leq \operatorname{ev}_9(u)
 \leq 9\frac{9^\ell-1}{8}.
\]
El mínimo del intervalo siguiente es una unidad mayor que el máximo del
anterior. Dentro de cada intervalo, la división de $n-1$ por nueve determina
de manera única el dígito
\[
 d=(n-1)\bmod9+1
\]
y el cociente que se codifica recursivamente. El proceso termina porque el
cociente decrece estrictamente.
\end{proof}

Ejemplos:
\[
 9\longleftrightarrow(9),\qquad
 10\longleftrightarrow(1,1),\qquad
 81\longleftrightarrow(9,8),\qquad
 1000\longleftrightarrow(1,3,3,1).
\]

\section{Fibras de las proyecciones residuales en los levantamientos enteros}

Para este bloque distinguimos dos alfabetos. La forma normal
\(\Wnine\) usa los dígitos \(\{1,\ldots,9\}\), escritos desde el menos
significativo. En cambio, \(\mathscr B_9\) designa las escrituras posicionales
convencionales en base nueve, con dígitos \(\{0,\ldots,8\}\) y el más
significativo a la izquierda. No se identifican ambos tipos.

Sea
\[
 \mathscr W_{10}^{(2)}=\{0,\ldots,9\}^2,\qquad
 \rho_{10}(a,b)=(b,a),\qquad
 \operatorname{ev}_{10}(a,b)=10a+b,
\]
y sea \(q_9:\mathbb N\to\mathbb Z/9\mathbb Z\) la reducción residual.

\begin{lemma}[Descenso residual de la reversión decimal]
\label{int:pf84:SYN30-RHO-DESC}
La reversión sobre palabras decimales de longitud dos es involutiva y su
sombra módulo nueve es la identidad:
\[
 \rho_{10}^2=I,\qquad
 q_9\!\left(\operatorname{ev}_{10}(\rho_{10}(a,b))\right)
 =
 q_9\!\left(\operatorname{ev}_{10}(a,b)\right).
\]
La longitud fija conserva los ceros iniciales.
\end{lemma}

\begin{proof}
La doble permutación recupera \((a,b)\). Como \(10\equiv1\pmod9\),
\[
 10b+a\equiv a+b\equiv10a+b\pmod9.
\]
Una palabra cuya doble reversión no fuese ella misma o cuya suma de cifras
cambiase módulo nueve refutaría el lema.
\end{proof}

\begin{lemma}[Descenso residual del cuadrado]
\label{int:pf84:SYN30-SQ-DESC}
La aplicación \(s(n)=n^2\) induce el endomapa residual
\[
 [n]_9\longmapsto[n]_9^2,
 \qquad q_9(s(n))=q_9(n)^2.
\]
\end{lemma}

\begin{proof}
La compatibilidad de una congruencia con el producto da
\(n\equiv m\pmod9\Rightarrow n^2\equiv m^2\pmod9\). Un entero para el que
\(q_9(n^2)\ne q_9(n)^2\) refutaría el enunciado.
\end{proof}

\begin{proposition}[Fibras residuales de la reversión entera]
\label{int:pf84:SYN30-RHO-NONFAITH}
La sombra residual no reconstruye el levantamiento entero
\(\widehat\rho_{10}=\operatorname{ev}_{10}\circ\rho_{10}\). En efecto,
\[
 18\equiv27\pmod9,\qquad
 \widehat\rho_{10}(1,8)=81\ne72=\widehat\rho_{10}(2,7).
\]
\end{proposition}

\begin{proof}
Ambos enteros pertenecen a la clase residual cero, pero sus reversiones son
distintas. Por tanto no existe
\(\overline\rho:\mathbb Z/9\mathbb Z\to\mathbb N\) que factorice
simultáneamente esos dos valores.
\end{proof}

\begin{proposition}[Fibras residuales del cuadrado entero]
\label{int:pf84:SYN30-SQ-NONFAITH}
La misma sombra tampoco reconstruye el cuadrado entero:
\[
 18\equiv27\pmod9,\qquad
 18^2=324\ne729=27^2.
\]
\end{proposition}

\begin{proof}
Una función \(\overline s:\mathbb Z/9\mathbb Z\to\mathbb N\) que
factorizase \(s\) debería asignar a la misma clase cero los dos valores
\(324\) y \(729\), lo cual es imposible.
\end{proof}

\begin{corollary}[Testigo cruzado \(27\)--\(72\)--\(729\)]
\label{int:pf84:SYN30-CROSS-WITNESS}
\[
 \widehat\rho_{10}(2,7)=72,\qquad
 27^2=729=9^3=3^6,
\]
y, en el alfabeto convencional \(\mathscr B_9\),
\[
 27=30_9,\qquad72=80_9,\qquad729=1000_9.
\]
Los tres enteros tienen sombra residual cero y raíz digital positiva nueve.
\end{corollary}

\begin{proof}
Las dos operaciones se calculan directamente; las divisiones sucesivas por
nueve producen las tres escrituras. El resultado fallaría si se confundiese
\(\mathscr B_9\) con \(\Wnine\), si una de las escrituras fuese incorrecta
o si se afirmase que el descenso residual fiel reconstruye la operación y su
genealogía completa.
\end{proof}

\section{Sumador nativo y propagación del cociente}

\begin{algorithm}[Suma de palabras]
Se recorren dos palabras por posiciones, con acarreo inicial \(c_0=0\). Si en
una posición hay dos dígitos, se aplica \eqref{int:eq:local-plus}; si sólo hay uno,
ese dígito se toma como residuo provisional y el cociente provisional es cero.
Cuando \(c_j\ne0\), se aplica una segunda celda aditiva al residuo provisional
y a \(c_j\). Se publica el nuevo residuo y se transporta la suma de los
cocientes a la posición siguiente. Si al terminar queda un acarreo no nulo, se
publica como último dígito. De este modo las celdas locales nunca reciben el
símbolo ausente \(0\), que no pertenece a \(\Dnine\).
\end{algorithm}

\begin{lemma}[Invariante del sumador]
Si $u\boxplus_\#v$ es la palabra producida por el algoritmo anterior,
\[
 \operatorname{ev}_9(u\boxplus_\#v)
 =\operatorname{ev}_9(u)+\operatorname{ev}_9(v).
\]
El acarreo intermedio está en $\{0,1,2\}$.
\end{lemma}

\begin{proof}
Tras procesar las posiciones $0,\ldots,k-1$, sea $c_k$ el cociente pendiente.
Las identidades locales implican
\[
 \sum_{j<k}(u_j+v_j)9^j
 =\sum_{j<k}w_j9^j+c_k9^k,
\]
incluyendo como cero los dígitos ausentes. Cada nueva celda conserva esta
igualdad. El máximo local es $9+9+2=20$, de modo que $c_{k+1}\leq2$. Al
terminar se publica el último cociente, y queda la identidad global.
\end{proof}

\section{Multiplicación por un dígito}

\begin{algorithm}[Transductor escalar]
Fijado $d\in\Dnine$, cada dígito $u_j$ produce
\[
 u_jd=9q_j+r_j,
 \qquad
 (r_j,q_j)=\bigl(R^\times(u_j,d),q^\times(u_j,d)\bigr).
\]
Sea $c_j\in\{0,\ldots,9\}$ el cociente entrante. Como $0\notin\Dnine$,
la transición se define por ramas:
\[
 (w_j,c_{j+1})=
 \begin{cases}
 (r_j,q_j),&c_j=0,\\[1mm]
 \bigl(R^+(r_j,c_j),q_j+q^+(r_j,c_j)\bigr),&c_j\in\Dnine.
 \end{cases}
\]
Así, ninguna aplicación local se evalúa fuera de su dominio. La primera
rama expresa la ausencia de cociente pendiente; la segunda normaliza la
colisión mediante la aplicación aditiva.
\end{algorithm}

\begin{lemma}[Invariante escalar]
La palabra $d\odot_\#u$ producida por el transductor satisface
\[
 \operatorname{ev}_9(d\odot_\#u)
 =d\operatorname{ev}_9(u).
\]
El cociente multiplicativo pertenece a $\{0,\ldots,9\}$.
\end{lemma}

\begin{proof}
El argumento es el del sumador, sustituyendo cada sumando local por
$u_jd$. Si $c_j=0$, la identidad
$u_jd=9q_j+r_j$ conserva directamente la igualdad parcial. Si
$c_j\in\Dnine$, la segunda división euclídea da
\[
 r_j+c_j=9q^+(r_j,c_j)+R^+(r_j,c_j),
\]
y conserva la misma igualdad después de transportar
$q_j+q^+(r_j,c_j)$. El máximo $9\cdot9+9=90$ y la forma de resto en
$\Dnine$ mantienen el cociente en la cota declarada. El último cociente se
publica como dígito.
\end{proof}

\section{Producto global por Horner}

\begin{definition}[Producto nativo]
Sea $v=(v_0,\ldots,v_{m-1})$. Partiendo de $a_m=\varnothing$, se define para
$j=m-1,\ldots,0$
\[
 a_j=(9\odot_\#a_{j+1})\boxplus_\#(v_j\odot_\#u).
\]
El resultado es
\[
 u\boxtimes_\#v:=a_0.
\]
\end{definition}

Esta definición no contiene
\[
 \operatorname{dec}_9(\operatorname{ev}_9(u)\operatorname{ev}_9(v)).
\]
La evaluación sólo aparecerá en la demostración del entrelazamiento.

\begin{theorem}[Multiplicatividad nativa]
Para todas las palabras $u,v\in\Wnine$,
\begin{equation}
 \label{int:eq:ev-multiplicative}
 \boxed{
 \operatorname{ev}_9(u\boxtimes_\#v)
 =\operatorname{ev}_9(u)\operatorname{ev}_9(v).}
\end{equation}
La palabra vacía es absorbente y $(1)$ es la unidad.
\end{theorem}

\begin{proof}
Por los dos lemas anteriores,
\[
 \operatorname{ev}_9(a_j)
 =9\operatorname{ev}_9(a_{j+1})
  +v_j\operatorname{ev}_9(u).
\]
La iteración de Horner da
\[
 \operatorname{ev}_9(a_0)
 =\operatorname{ev}_9(u)
  \sum_{j=0}^{m-1}v_j9^j.
\]
El último factor es $\operatorname{ev}_9(v)$. La unicidad de la forma normal
identifica la palabra resultante.
\end{proof}

\begin{corollary}[Monoide multiplicativo de palabras]
La estructura
\[
\WordMonoid=(\Wnine,\boxtimes_\#,(1))
\]
es un monoide conmutativo; la palabra vacía es su elemento absorbente. Su
construcción no utiliza la multiplicación global de enteros.
\end{corollary}

El ejemplo de frontera más corto es
\[
 (9)\boxtimes_\#(9)=(9,8),
 \qquad 9+8\cdot9=81.
\]
Las \(17\) celdas locales de la multiplicación producen
\[
 (9,8)\boxtimes_\#(9,8)=(9,8,8,8),
 \qquad \operatorname{ev}_9(9,8,8,8)=6561.
\]

\section{Linealización matricial del estado genealógico y construcción del entrelazador completamente positivo}

\begin{definition}[Sectores de Hilbert]
Sean
\[
 \Krad=\ell^2(\Wnine^+),
 \qquad
 \Hplus=\ell^2(\N_{\geq1}).
\]
Definimos sobre las bases canónicas
\[
 W_\times e_u=e_{\operatorname{ev}_9(u)}.
\]
En \(\Hplus\), el operador número se denota por
\[
 \mathsf N e_n=n e_n,
 \qquad \Spec(\mathsf N)=\N_{\geq1}.
\]
El vacío puede añadirse separadamente como un vector de vacío $e_0$; no se
confunde con $\Hplus$, cuyo espectro comienza en uno.
\end{definition}

\begin{theorem}[Unitariedad y cuadrado multiplicativo]
$W_\times$ es unitario. En el núcleo algebraico
$\mathbb C[\Wnine^+\times\Wnine^+]$, sea
\[
 \mu_\#(e_u\otimes e_v)=e_{u\boxtimes_\#v},
 \qquad
 \mu_{\N}(e_a\otimes e_b)=e_{ab}.
\]
Entonces vale el cuadrado
\[
\begin{CD}
\Krad\otimes_{\mathrm{alg}}\Krad @>{W_\times\otimes W_\times}>>
\Hplus\otimes_{\mathrm{alg}}\Hplus\\
@V{\mu_\#}VV @VV{\mu_{\N}}V\\
\Krad @>>{W_\times}> \Hplus
\end{CD}
\]
y, por tanto, el cuadrado conmuta.
\end{theorem}

\begin{proof}
La biyección de formas normales envía una base ortonormal sobre otra y prueba
la unitariedad. En un tensor básico, ambos caminos producen
$e_{\operatorname{ev}_9(u)\operatorname{ev}_9(v)}$ por
\eqref{int:eq:ev-multiplicative}. La linealidad da la igualdad en el núcleo.
\end{proof}

Los mapas $\mu_\#$ y $\mu_{\N}$, definidos inicialmente sobre los respectivos
núcleos algebraicos, son densamente definidos y no acotados. Para
\(c=(c_{u,v})\in\ell^2(\Wnine^+\times\Wnine^+)\), cada suma sobre una fibra de producto
es finita. El adjunto de \(\mu_\#\) contiene todos los vectores de soporte
finito en su dominio, de modo que \(\mu_\#\) es clausurable. El dominio de su
clausura es
\[
 \Dom(\overline{\mu_\#})=
 \left\{c\in\ell^2(\Wnine^+\times\Wnine^+):
 \sum_{w\in\Wnine^+}
 \left|\sum_{u\boxtimes_\#v=w}c_{u,v}\right|^2<\infty
 \right\}.
\]
Cada fibra es finita porque $\operatorname{ev}_9$ la identifica con los pares
ordenados de divisores de un entero. La misma descripción vale para
\(\overline{\mu_{\N}}\) después de aplicar \(W_\times\otimes W_\times\), y
$W_\times$ entrelaza los grafos de ambas clausuras.

\section{Coproducto y selector nativos}

La apertura multiplicativa se define ahora antes de mirar el codominio:
\[
 \Delta_\#e_w
 =\sum_{u\boxtimes_\#v=w}e_u\otimes e_v.
\]
La suma se obtiene enumerando productos del transductor. No se pregunta si
un entero ordinario divide a otro.
En el sector entero definimos, sobre el núcleo algebraico correspondiente,
\[
 \Delta_\times e_n
 =\sum_{\substack{a,b\in\N_{\geq1}\\ab=n}}e_a\otimes e_b.
\]
Ambos coproductos enumeran pares ordenados.

\begin{theorem}[Coasociatividad nativa]
En el núcleo algebraico,
\[
 (\Delta_\#\otimes I)\Delta_\#
 =(I\otimes\Delta_\#)\Delta_\#.
\]
Además,
\[
 (W_\times\otimes W_\times)\Delta_\#
 =\Delta_\times W_\times.
\]
\end{theorem}

\begin{proof}
Ambos lados de la primera igualdad enumeran una vez cada triple nativo
$a\boxtimes_\#b\boxtimes_\#c=w$. La segunda igualdad sigue de la biyección
entre fibras que induce \eqref{int:eq:ev-multiplicative}.
\end{proof}

\begin{theorem}[Coproducto aritmético completo]
\label{int:pf84:C06}
En el núcleo algebraico
\(c_{00}(\mathbb N_{\geq1})\subset\ell^2(\mathbb N_{\geq1})\), sea
\[
 \Delta_\times e_n=\sum_{ab=n}e_a\otimes e_b.
\]
Su clausura tiene dominio
\[
 \operatorname{Dom}(\overline{\Delta_\times})
 =
 \left\{x=(x_n):
 \sum_{n\geq1}d(n)|x_n|^2<\infty\right\},
\]
es coasociativa en el núcleo y, para \(Qe_n=ne_n\), satisface
\[
 (\log Q\otimes I+I\otimes\log Q)\Delta_\times
 =\Delta_\times\log Q.
\]
\end{theorem}

\begin{proof}
Los soportes de \(\Delta_\times e_n\) son ortogonales para enteros distintos
y su norma cuadrada es el número \(d(n)\) de divisores, lo que da el dominio
ponderado de la clausura. Las dos coasociaciones reindexan la misma suma
sobre \(abc=n\). Finalmente, \(\log a+\log b=\log n\) en cada término
prueba el entrelazamiento. Una triple factorización contada de forma distinta
o el fallo de esta identidad refutaría el teorema.
\end{proof}

\begin{theorem}[Coproducto reducido y selector de irreducibles]
\label{int:pf84:C07}
En el sector entero, el coproducto reducido
\[
 \Delta_\times^{\rm red}e_n
 =\sum_{\substack{ab=n\\a,b>1}}e_a\otimes e_b
\]
es clausurable. El operador
\[
 N_\times^{\rm red}
 =\bigl(\overline{\Delta_\times^{\rm red}}\bigr)^*
  \overline{\Delta_\times^{\rm red}}
\]
es positivo autoadjunto y diagonal:
\[
 N_\times^{\rm red}e_n=d_{\rm red}(n)e_n.
\]
Su núcleo está generado por \(e_1\) y por los \(e_p\) con \(p\) irreducible;
por tanto,
\[
 P_{\Irr}
 =\mathbf1_{\{0\}}(N_\times^{\rm red})
  -\lvert e_1\rangle\langle e_1\rvert
\]
selecciona los irreducibles y excluye la unidad.

En el sector nativo, el coproducto correspondiente es
\[
 \widetilde\Delta_\#e_w
 =\sum_{\substack{u\boxtimes_\#v=w\\u,v\ne(1)}}e_u\otimes e_v.
\]
El entrelazador \(W_\times\) transporta
\(\Delta_\times^{\rm red}\), \(N_\times^{\rm red}\) y \(P_{\Irr}\) hacia
\(\widetilde\Delta_\#\),
\[
 A_\#:=
 \bigl(\overline{\widetilde\Delta_\#}\bigr)^*
 \overline{\widetilde\Delta_\#},
 \qquad
 P_{\Irr,\#}
 =(I-P_{(1)})\mathbf1_{\{0\}}(A_\#),
 \qquad
 P_{(1)}=\lvert e_{(1)}\rangle\langle e_{(1)}\rvert.
\]
\end{theorem}

\begin{proof}
Los soportes de las imágenes de dos vectores base distintos son ortogonales.
La norma cuadrada de
\(\Delta_\times^{\rm red}e_n\) es \(d_{\rm red}(n)\), que se anula
exactamente para \(n=1\) o \(n\) irreducible. Esto prueba la diagonalización
y el núcleo declarado; restar el proyector de \(e_1\) elimina la unidad.
La multiplicatividad de \(W_\times\) identifica las fibras enteras con las
fibras nativas y transporta los tres operadores. Un compuesto en el núcleo,
un irreducible fuera de él o la inclusión de la unidad en \(P_{\Irr}\)
refutarían el teorema.
\end{proof}

\begin{proposition}[Coincidencia nonádica de primos y compuestos]
\label{int:pf84:B05}
Sea \(\gcd(n,90)=1\),
\(\tau_n=\operatorname{ord}_n(10)\) y
\[
 M_n=\frac{10^{\tau_n}-1}{n}.
\]
Entonces \(M_n\equiv0\pmod9\), tanto para primos como para compuestos. En
particular, el cierre nonádico y el período decimal aislados no son un
selector de irreducibles.
\end{proposition}

\begin{proof}
Como \(9\mid10^{\tau_n}-1\) y \(n\) es invertible módulo nueve, la
divisibilidad por nueve sobrevive al cociente. Los primos \(7,13\) y los
compuestos \(77,91,143\) comparten el período seis; los números de Carmichael
aportan controles adicionales. El enunciado quedaría refutado sólo si la
congruencia o el período separasen todos los primos de todos los compuestos
en el dominio declarado.
\end{proof}

Esta afirmación es estrictamente aritmética: identifica fibras de
factorización y los operadores que éstas inducen sobre el sector de formas
normales.

\ifdefined\HMTInternalTraceability
\section{Extensión monoidal del producto nativo a historias compatibles}
\label{int:pf84:SYN-05D-HIST}

El resultado distingue el espacio \(\Krad\) de formas normales canónicas y
el espacio \(\mathcal K_{\#,\mathrm{hist}}\) de historias completas de
celdas y cocientes propagados.
Muchas historias pueden producir la misma forma normal. Por tanto,
$W_\times$ es unitario sobre $\Krad$, pero la proyección de una historia sobre
su valor impide que sea inyectivo sobre todas las clases de transporte.

Sea \(\mathscr X_n^{\rm hist}\) el conjunto de historias enriquecidas de
profundidad \(n\) y sea
\(\mathcal H_n=\ell^2(\mathscr X_n^{\rm hist})\). Denotemos por \(P_n\) la
proyección ortogonal sobre el subespacio generado por las historias
supervivientes y supongamos definida una multiplicación \(\mu_n\) sobre un
dominio algebraico de
\(\mathcal H_n\otimes_{\mathrm{alg}}\mathcal H_n\). Una extensión monoidal
del producto de palabras al sistema inverso consiste en una familia de
aplicaciones \(\iota_n:\Krad\to\mathcal H_n\) que conmuta con las
restricciones de profundidad y satisface, para cada \(n\),
\[
 P_n\,\iota_n\mu_\#
 =P_n\,\mu_n
 (\iota_n\otimes\iota_n)
\]
en el núcleo algebraico. Esta igualdad es la condición de compatibilidad que
define la extensión. Los teoremas precedentes no establecen la existencia ni
la unicidad de una familia \((\iota_n)_n\); demuestran el producto en el
monoide \(\WordMonoid\) y determinan la ecuación que una extensión al espacio
de historias debe satisfacer. La transferencia conserva, por tanto, el
estatuto de \emph{programa abierto}: no hereda la exactitud del producto en
\(\WordMonoid\). La inexistencia de una familia compatible, el fallo de
conmutación con los truncamientos, la pérdida de supervivencia, cociente,
acarreo, orientación o memoria, o el incumplimiento de la identidad anterior
son sus falsadores.
\fi

\section{Enumeración exhaustiva y necesidad de las coordenadas estructurales}

La enumeración finita comprende:
\[
\begin{array}{@{}lr@{}}
\toprule
\text{Prueba}&\text{Resultado}\\\midrule
\text{Celdas por aplicación local}&81\\
\text{Recorridos reversibles de numeración}&59050\\
\text{Pares del entrelazador}&532900\\
\text{Errores del entrelazador}&0\\
\text{Triples asociativos}&531441\\
\text{Errores asociativos}&0\\
\text{Ventana de coproducto}&1\text{ a }500\\
\text{Irreducibles obtenidos}&95\\
\bottomrule
\end{array}
\]

\begin{definition}[Cociente multiplicativo]
Conservar sólo $R^\times$ transforma $2\cdot5=10$ en el valor visible $1$.
Perturbar $q^\times(9,9)$ de $8$ a $7$ transforma $81$ en $72$. El cociente
local es necesario.
\end{definition}

\begin{definition}[Aplicación aditiva]
La aplicación multiplicativa produce productos parciales, pero no resuelve sus colisiones.
Eliminar $q^+$ rompe la normalización global; por ejemplo, deja de reproducir
los cocientes intermedios de $10\cdot18=180$.
\end{definition}

\begin{definition}[Coordenadas y símbolo de borde]
El cero ordinario no es un dígito de la construcción. La palabra vacía es
$0$ y el símbolo $9$ conserva su valor y su cociente. Mezclar las coordenadas
posicionales indexadas desde cero con las etiquetas $1,\ldots,9$ modifica las
celdas.
\end{definition}

\begin{proposition}[Alcance exacto de las ablaciones y del censo finito]
\label{int:pf84:D10}
Las ablaciones anteriores prueban, dentro de las convenciones declaradas:
\begin{enumerate}
  \item sin \(q^\times\), \(2\cdot5=10\) colapsa al residuo visible \(1\);
  \item cambiar \(q^\times(9,9)=8\) por \(7\) transforma \(81\) en \(72\);
  \item sin \(q^+\), falla la normalización de las colisiones aditivas.
\end{enumerate}
La enumeración exhaustiva comprende las \(2\cdot81=162\) celdas de las dos
aplicaciones, los \(729\) triples locales, \(532\,900\) pares de palabras y
los \(95\) irreducibles no mayores que \(500\). En estos dominios finitos se
demuestran las identidades indicadas, pero no se prueba la supervivencia
coinductiva de toda historia TPK.
\end{proposition}

\begin{proof}
Los tres contraejemplos se calculan directamente con
\eqref{int:eq:local-plus} y \eqref{int:eq:local-times}. La enumeración de los cuatro
dominios finitos produce los cardinales escritos en el enunciado y no añade
ninguna identidad coinductiva a las que ya se han demostrado en el monoide
de palabras.
\end{proof}

\section{Estructura algebraica del monoide de palabras}

La construcción demuestra, en el monoide \(\WordMonoid\):
\begin{enumerate}
 \item una forma normal única para cada entero no negativo;
 \item suma y producto globales construidos sólo con las dos aplicaciones locales;
 \item conservación simbólica del valor bajo la propagación del cociente;
 \item un unitario $W_\times$ y el cuadrado multiplicativo exacto;
 \item un coproducto divisor y un selector irreducible definidos nativamente;
 \item la necesidad de conservar la clase de transporte cuando se abandona el
 sector de formas normales.
\end{enumerate}

El paso al espacio enriquecido se reduce ahora a un problema de existencia:
construir una familia \((\iota_n)_n\) compatible con las restricciones de
profundidad y demostrar que satisface la ecuación de la sección anterior.
Esta cuestión es distinta del entrelazador exacto sobre
\(\WordMonoid\), que ya queda determinado por la forma normal y las
aplicaciones locales.
